% Worked solutions for Listas/lista-r2b-bond-valuation.md (15 exercises).
% Compile: pdflatex twice; check: pdffonts (Type 1 only, uni = yes).
% Every number is asserted in lista-r2b-bond-valuation-solutions_check.py.
\documentclass[11pt,a4paper]{article}
\usepackage[utf8]{inputenc}
\usepackage[T1]{fontenc}
\usepackage{lmodern}   % Type 1 vector fonts (replaces PK bitmaps)
\usepackage{cmap}      % ToUnicode CMap -> copyable text and accents
\usepackage[english]{babel}
\usepackage{amsmath,amssymb,amsfonts}
\usepackage[a4paper,margin=2.2cm]{geometry}
\usepackage{booktabs}
\usepackage{enumitem}
\usepackage{fancyhdr}
\usepackage{xcolor}
\usepackage{microtype}

\pagestyle{fancy}\fancyhf{}
\lhead{\footnotesize Finan\c{c}as II --- Corporativas (MPE Insper)}
\rhead{\footnotesize List R2b (Bond valuation) --- Worked solutions}
\cfoot{\footnotesize\thepage}
\renewcommand{\headrulewidth}{0.3pt}
\setlength{\headheight}{14pt}
\setlist{topsep=2pt,itemsep=1pt}

\newcommand{\exsec}[1]{\section*{\color{blue!55!black}#1}\addcontentsline{toc}{section}{#1}}
\newcommand{\notation}{\par\noindent\textbf{Notation.}}
\newcommand{\step}[2]{\par\medskip\noindent\textbf{Step #1 --- #2.}~}
\newcommand{\intuition}{\par\medskip\noindent\textbf{Intuition.}~}
\newcommand{\bookans}{\par\medskip\noindent\textbf{Book answer.}~}
\newcommand{\refbib}{\par\smallskip\noindent\textbf{Ref:}~}
\newcommand{\asked}{\par\medskip\noindent\textbf{What is asked.}~}
\newcommand{\approach}{\par\smallskip\noindent\textbf{Approach.}~}
\newcommand{\answer}[1]{\par\smallskip\noindent\fbox{\textbf{Answer #1.}}~}
\newcommand{\pitfalls}{\par\medskip\noindent\textbf{Common mistakes.}}
\newcommand{\E}{\mathbb{E}}

\usepackage[hyphens]{url}
%% STATEMENT-PREAMBLE BEGIN (tools/statements_into_solutions.py)
\usepackage{longtable,array,calc}
\providecommand{\tightlist}{\setlength{\itemsep}{0pt}\setlength{\parskip}{0pt}}
\providecommand{\href}[2]{#2}
\providecommand{\real}[1]{#1}
\makeatletter\@ifundefined{c@none}{\newcounter{none}}{}\makeatother  % pandoc's \LTcaptype{none}
\newcommand{\stmtbegin}{\par\medskip\noindent\textbf{Statement.}\par\nopagebreak\begingroup\small\setlength{\tabcolsep}{3pt}\sloppy}
\newcommand{\stmtend}{\par\endgroup\nopagebreak\noindent\rule{\linewidth}{0.4pt}\par\medskip}
\DeclareUnicodeCharacter{2192}{\ensuremath{\rightarrow}}
\DeclareUnicodeCharacter{2212}{\ensuremath{-}}
\DeclareUnicodeCharacter{2022}{\textbullet{}}
\DeclareUnicodeCharacter{2248}{\ensuremath{\approx}}
\DeclareUnicodeCharacter{2032}{\ensuremath{^{\prime}}}
\DeclareUnicodeCharacter{2081}{\ensuremath{_{1}}}
\DeclareUnicodeCharacter{03B8}{\ensuremath{\theta}}
\DeclareUnicodeCharacter{03BB}{\ensuremath{\lambda}}
\DeclareUnicodeCharacter{03C8}{\ensuremath{\psi}}
\DeclareUnicodeCharacter{03B1}{\ensuremath{\alpha}}
\DeclareUnicodeCharacter{25A0}{\rule{0.7em}{0.7em}}
%% STATEMENT-PREAMBLE END
\begin{document}
\begin{center}{\Large\bfseries Exercise list R2b --- Bond valuation, the term structure and the cost of debt}\\[4pt]{\large Worked solutions}\end{center}

\noindent Statements are in \url{Listas/lista-r2b-bond-valuation.md}; each exercise is cited by number and printed page and its full text is reproduced right before its solution. Short answers: \url{lista-r2b-bond-valuation-solutions-short.pdf}. Editions are the local copies: Berk = Berk, DeMarzo \& Harford, \emph{Fundamentals of Corporate Finance}, 2nd ed. (2011); Ross = Ross, Westerfield, Jordan \& Lamb, \emph{Fundamentos de Administra\c{c}\~ao Financeira}, 9th ed. (2013); Higgins = Koski, Higgins \& Mitton, \emph{Analysis for Financial Management}, 13th ed.; Grinblatt--Titman, \emph{Financial Markets and Corporate Strategy}, 2nd ed. (2002). Ross writes decimals with commas and \$1.000 for one thousand; these solutions use points. Numbers labelled ``illustration'' are not part of the statement.

\medskip\noindent\textbf{Formulas used throughout.} With coupon $C$ per period, face value $F$, $n$ periods and yield $y$ per period,
\begin{align}
  P &= C\cdot A(y,n) + \frac{F}{(1+y)^n}, \qquad A(y,n) = \frac{1-(1+y)^{-n}}{y} \quad\text{(annuity factor)}, \label{r2be0-bond}\\
  P &= \sum_{t=1}^{n}\frac{CF_t}{(1+r_t)^t} \quad\text{(spot rates $r_t$, Berk Eq.~6.4)}. \label{r2be0-spot}
\end{align}
A semiannual bond quoted at an APR $y^{\text{APR}}$ has $y=y^{\text{APR}}/2$ per period and $C=$ coupon rate$\times F/2$.

\tableofcontents

% =============================================================================
\exsec{1. Zero-coupon prices and the yield curve --- Berk 6.4, p.~172}
%% BEGIN-STATEMENT 1
\stmtbegin
\textbf{Source:} Berk, Ch. 6, Exercise 6.4, p.~172 (PDF 203) · Answer in book: no

The following table summarizes prices of various default-free zero-coupon bonds (expressed as a
percentage of face value):

{\def\LTcaptype{none} % do not increment counter
\begin{longtable}[]{@{}llllll@{}}
\toprule\noalign{}
Maturity (years) & 1 & 2 & 3 & 4 & 5 \\
\midrule\noalign{}
\endhead
\bottomrule\noalign{}
\endlastfoot
Price (per \$100 face value) & \$95.51 & \$91.05 & \$86.38 & \$81.65 & \$76.51 \\
\end{longtable}
}

\begin{enumerate}
\def\labelenumi{\alph{enumi}.}
\item
  Compute the yield to maturity for each bond.
\item
  Plot the zero-coupon yield curve (for the first five years).
\item
  Is the yield curve upward sloping, downward sloping, or flat?
\end{enumerate}
\stmtend
%% END-STATEMENT

\asked (a) The yield to maturity of each of the five default-free zeros from its price per \$100 of face value.
(b) A plot of yield against maturity. (c) Whether the curve slopes up, down or is flat.

\approach A zero-coupon bond has a single cash flow, the face value at maturity, so its price is that face
value discounted at its own yield: $P_n = 100/(1+YTM_n)^n$. Inverting gives the yield. The zero-coupon
yields are the \emph{spot rates} of block 2A.1, the discount rates used to price any other default-free bond.

\notation
\begin{itemize}
  \item $n\in\{1,\dots,5\}$: maturity in years; $P_n$: price per \$100 face value; $YTM_n$: annual yield (spot rate) for maturity $n$.
\end{itemize}

\step{1}{Invert the pricing formula} The price of a zero is
\begin{align}
  P_n = \frac{100}{(1+YTM_n)^n} \;\Longrightarrow\; (1+YTM_n)^n = \frac{100}{P_n} \;\Longrightarrow\; YTM_n = \left(\frac{100}{P_n}\right)^{1/n}-1. \label{r2be1-inv}
\end{align}

\step{2}{Substitute each price} We compute the ratio $100/P_n$, take the $n$th root and subtract one.
\begin{center}
\begin{tabular}{@{}cccc c@{}}\toprule
$n$ & $P_n$ & $100/P_n$ & $(100/P_n)^{1/n}$ & $YTM_n$\\\midrule
1 & 95.51 & 1.047011 & $1.047011^{1}=1.047011$ & 4.70\%\\
2 & 91.05 & 1.098298 & $1.098298^{1/2}=1.047997$ & 4.80\%\\
3 & 86.38 & 1.157675 & $1.157675^{1/3}=1.050015$ & 5.00\%\\
4 & 81.65 & 1.224740 & $1.224740^{1/4}=1.051988$ & 5.20\%\\
5 & 76.51 & 1.307019 & $1.307019^{1/5}=1.055009$ & 5.50\%\\\bottomrule
\end{tabular}
\end{center}
\answer{(a)} $YTM_1=4.70\%$, $YTM_2=4.80\%$, $YTM_3=5.00\%$, $YTM_4=5.20\%$, $YTM_5=5.50\%$.

\step{3}{Plot} The curve joins the points $(1,4.70)$, $(2,4.80)$, $(3,5.00)$, $(4,5.20)$, $(5,5.50)$ (maturity in years, yield in \%).
\answer{(b)} A rising line through those five points, steepening after year 3.

\step{4}{Slope} Each yield exceeds the previous one: $4.70<4.80<5.00<5.20<5.50$.
\answer{(c)} Upward sloping.

\intuition Longer zeros are cheaper per year of waiting than a constant rate would imply: the market
asks a higher annual rate to lend for longer. These five numbers are all one needs to price any default-free
bond paying on those dates (Exercise~8).
\pitfalls
\begin{itemize}
  \item Computing $(100-P_n)/P_n$ and calling it the annual yield: that is the \emph{total} return over $n$ years, not an annual rate.
  \item Dividing the total return by $n$ (simple interest) instead of taking the $n$th root.
\end{itemize}
\refbib{Berk \S6.2 ``Zero-Coupon Bonds'', pp.~147--151.}

% =============================================================================
\exsec{2. Holding-period return on a bond --- Higgins 5.5, p.~188}
%% BEGIN-STATEMENT 2
\stmtbegin
\textbf{Source:} Higgins (Koski--Higgins--Mitton), Ch. 5, Exercise 5.5, p.~188 (PDF 209) · Answer in
book: yes, p.~410

The return an investor earns on a bond over a period of time is known as the holding period return,
defined as interest income plus or minus the change in the bond's price, all divided by the
beginning bond price.

\begin{enumerate}
\def\labelenumi{\alph{enumi}.}
\item
  What is the holding period return on a bond with a par value of \$1,000 and a coupon rate of 6
  percent if its price at the beginning of the year was \$1,050 and its price at the end was \$940?
  Assume interest is paid annually.
\item
  Can you give two reasons the price of the bond might have fallen over the year?
\end{enumerate}
\stmtend
%% END-STATEMENT

\asked (a) The return over one year on a bond bought at \$1,050 and worth \$940 a year later, after one
annual coupon at 6\% of \$1,000. (b) Two reasons why its price could have fallen.

\approach The holding-period return is the definition given in the statement: income plus capital gain
over the starting price. Part (b) uses the pricing formula \eqref{r2be0-bond}: the price falls if the
yield the market demands rises, either because riskless rates rise or because the issuer's credit risk rises.

\notation
\begin{itemize}
  \item $P_0=1{,}050$, $P_1=940$: prices at the start and end of the year; $I$: coupon received; $HPR$: holding-period return.
\end{itemize}

\step{1}{Coupon} The annual coupon is
\begin{align}
  I = 0.06\times 1{,}000 = 60. \label{r2be2-I}
\end{align}

\step{2}{Return} We add the coupon to the price change and divide by the starting price:
\begin{align}
  HPR = \frac{I + (P_1-P_0)}{P_0} = \frac{60 + (940-1{,}050)}{1{,}050} = \frac{60-110}{1{,}050} = \frac{-50}{1{,}050} = -4.76\%. \label{r2be2-hpr}
\end{align}
\answer{(a)} $HPR=-4.76\%$.

\step{3}{Why the price fell} From \eqref{r2be0-bond}, with coupon and face fixed, $P$ falls only if the
discount rate $y$ rises (the passage of time alone pulls a premium bond toward par, but a 6\% bond at
\$1,050 would only lose a few dollars that way, not \$110).
\answer{(b)} (i) Market interest rates rose (interest-rate risk); (ii) the market judged the issuer riskier
(a downgrade or worse prospects raised the default premium). Either raises the yield and lowers the price.

\intuition The coupon is fixed, so all news reaches the bondholder through the price. Even a default-free
bond has a risky one-year return unless it matures exactly at the end of the holding period.
\pitfalls
\begin{itemize}
  \item Forgetting the coupon: $(940-1{,}050)/1{,}050=-10.5\%$ is the capital loss only.
  \item Dividing by the face value \$1,000 instead of the price paid.
\end{itemize}
\bookans Matches the book's answer, $-4.76$ percent, with interest rates or perceived risk rising as the reasons (p.~410).
\refbib{Higgins Ch.~5, ``Bonds'' and ``Bond Ratings'', pp.~143--150; Ch.~7 for bond pricing.}

% =============================================================================
\exsec{3. Price of a seasoned semiannual bond --- Ross 7.6, p.~232}
%% BEGIN-STATEMENT 3
\stmtbegin
\textbf{Source:} Ross, Cap. 7, Exercise 7.6, p.~232 (PDF 256) · Answer in book: no

6. Preços de títulos de dívida {[}OA2{]} A Companhia Pereira emitiu títulos de 11 anos há um ano,
com taxa de cupom de 6,9\%. Os pagamentos são semestrais. Se o YTM sobre esses títulos for de 7,4\%,
qual é o seu preço corrente?

\paragraph{Context needed to solve it}

\emph{From Ross §7.1 ``Valores e retornos dos títulos de dívida'', p.~195 (PDF 219) --- the bond
value formula:}

Com base em nossos exemplos, agora podemos escrever a expressão geral do valor de um título de
dívida. Se um título tem (1) um valor de face (F) pago no vencimento; (2) um cupom (C) pago por
período; (3) t períodos até o vencimento; e (4) um retorno de r por período, seu valor é:

\[\text{Valor do título de dívida} = C \times [1 - 1/(1+r)^t]/r + F/(1+r)^t\]

\emph{From Ross Example 7.1 ``Cupons semestrais'', p.~195 (PDF 219):}

Os títulos de dívida (bonds) emitidos nos Estados Unidos normalmente fazem pagamentos de cupom duas
vezes ao ano. Assim, se um bond comum tiver uma taxa de cupom de 14\%, o seu titular terá um total
de \$140 por ano, mas esses \$140 virão em dois pagamentos de \$70 cada um. {[}\ldots{]} Os retornos
desses títulos são cotados em taxa nominal anual; a taxa cotada é igual à taxa efetiva do período
multiplicada pelo número de períodos. Nesse caso, com um retorno cotado a 16\% e pagamentos
semestrais, o verdadeiro retorno é de 8\% a cada seis meses.

\emph{Face value:} the statement gives none; Ross's bond examples and problems use a face value of
\$1.000 (e.g.~Problems 7.1--7.2 on p.~229), and the answer is given per \$1.000 of face value.
\stmtend
%% END-STATEMENT

\asked The price today of a bond issued one year ago with an 11-year maturity, a 6.9\% coupon paid
twice a year, when its yield to maturity is 7.4\% (APR, semiannual).

\approach Ross's bond value formula \eqref{r2be0-bond} with half-year periods. Two conversions do the work:
the remaining life is $11-1=10$ years, i.e.\ 20 half-years, and both the coupon and the quoted yield are
halved. Face value \$1,000 (the book's convention; see the Context in the list).

\notation
\begin{itemize}
  \item $F=1{,}000$; $C=34.5$: coupon per half-year; $y=3.7\%$: yield per half-year; $n=20$ half-years.
\end{itemize}

\step{1}{Per-period inputs} We convert the annual figures to half-years:
\begin{align}
  C = \frac{0.069\times1{,}000}{2} = \frac{69}{2} = 34.5,\qquad y = \frac{0.074}{2} = 0.037,\qquad n = 2\times(11-1) = 20. \label{r2be3-in}
\end{align}

\step{2}{Factors} The compounding factor and the annuity factor are
\begin{align}
  (1.037)^{20} &= 2.068117,\qquad \frac{1}{(1.037)^{20}} = 0.483532, \label{r2be3-v}\\
  A(0.037,20) &= \frac{1-0.483532}{0.037} = \frac{0.516468}{0.037} = 13.958605. \label{r2be3-a}
\end{align}

\step{3}{Price} Substituting into \eqref{r2be0-bond}:
\begin{align}
  P = (34.5)(13.958605) + (1{,}000)(0.483532) = 481.57 + 483.53 = 965.10. \label{r2be3-P}
\end{align}
\answer{\unskip} $P=\$965.10$ per \$1,000 of face value (a discount bond: coupon 6.9\% $<$ yield 7.4\%).

\intuition The bond pays 6.9\% when the market wants 7.4\%, so it must sell below par; the discount
(\$34.90) is the present value of the coupon shortfall $(37-34.5)\times A(0.037,20)=2.5\times13.958605=34.90$.
\pitfalls
\begin{itemize}
  \item Using 11 periods or 22 half-years: the bond is one year old, so 10 years (20 half-years) remain.
  \item Discounting semiannual coupons at 7.4\% per period, or annual coupons at 3.7\%.
\end{itemize}
\refbib{Ross \S7.1 ``Valores e retornos dos t\'itulos de d\'ivida'', bond value formula and Example 7.1, pp.~192--196.}

% =============================================================================
\exsec{4. Callable vs.\ noncallable bonds --- Higgins 5.4, p.~188}
%% BEGIN-STATEMENT 4
\stmtbegin
\textbf{Source:} Higgins (Koski--Higgins--Mitton), Ch. 5, Exercise 5.4, p.~188 (PDF 209) · Answer in
book: no

Two 20-year bonds are identical in all respects except that one allows the issuer to call the bond
in return for \$1,000 cash at any time after five years, while the other contains no call
provisions. Will the yield to maturity on the two bonds differ? If so, which will be higher? Why?
\stmtend
%% END-STATEMENT

\asked Whether two otherwise identical 20-year bonds, one callable at \$1,000 after five years and one not,
have different yields; which is higher; and why.

\approach A call provision is an option held by the \emph{issuer}: the bondholder has sold it. By the law of
one price, the callable bond is a straight bond minus a call option, so its price is lower and, for the same
promised cash flows, its yield is higher.

\notation
\begin{itemize}
  \item $P^{\text{straight}}$, $P^{\text{callable}}$: prices; $V^{\text{call}}\ge 0$: value of the issuer's option to repay at \$1,000.
\end{itemize}

\step{1}{Decompose the callable bond} The holder of a callable bond owns the straight bond and is short the call:
\begin{align}
  P^{\text{callable}} = P^{\text{straight}} - V^{\text{call}} \le P^{\text{straight}}. \label{r2be4-dec}
\end{align}

\step{2}{When is the call exercised} The issuer calls when rates have fallen so that the straight bond is worth
more than \$1,000: it refinances at the lower rate. The bondholder loses exactly the gain she would otherwise
have had and must reinvest at the lower rate. The call caps her upside.

\step{3}{From price to yield} Both bonds promise the same coupons and face value if not called; a lower price for the
same promised flows means a higher discount rate solving \eqref{r2be0-bond}.
\answer{\unskip} Yes: the callable bond's yield to maturity is \emph{higher}, compensating investors for the option they give
the issuer (the call is exercised precisely when it hurts them: after a fall in rates).

\intuition Exercise~14 puts a number on $V^{\text{call}}$: about 12\% of face value for a Treasury callable within a year.
The yield difference is a price for an option, not a default premium.
\pitfalls
\begin{itemize}
  \item Thinking the callable bond is \emph{safer} because it may be repaid early; early repayment happens when the investor least wants it.
  \item Confusing the call (issuer's option) with a put (holder's option, which lowers the yield).
\end{itemize}
\refbib{Higgins Ch.~5, ``Bonds'' (call provisions), p.~145; Ross \S7.2 ``Op\c{c}\~ao de resgate antecipado'', p.~207.}

% =============================================================================
\exsec{5. Expected return vs.\ yield on a corporate bond --- Berk Critical Thinking 8, p.~172}
%% BEGIN-STATEMENT 5
\stmtbegin
\textbf{Source:} Berk, Ch. 6, Critical Thinking question 8, p.~172 (PDF 203) · Answer in book: no

Explain why the expected return of a corporate bond does not equal its yield to maturity.
\stmtend
%% END-STATEMENT

\asked Why the expected return of a corporate bond differs from its yield to maturity.

\approach Block 2A.10: the yield to maturity is the internal rate of return of the \emph{promised} cash flows,
the expected return is the rate that equates the price to the \emph{expected} cash flows. With default
risk the expected cash flows are smaller, so the same price implies a smaller rate.

\notation
\begin{itemize}
  \item $P$: price; $F$: promised payment (one-year zero, for clarity); $p$: default probability; $C$: recovery in default.
  \item $y$: promised yield, $P=F/(1+y)$; $r_D$: expected return (cost of debt), $P=[F(1-p)+Cp]/(1+r_D)$.
\end{itemize}

\step{1}{Two definitions, one price} From the two pricing equations,
\begin{align}
  1+y = \frac{F}{P}, \qquad 1+r_D = \frac{F(1-p)+Cp}{P}. \label{r2be5-defs}
\end{align}

\step{2}{Compare} Since $C<F$, $F(1-p)+Cp = F - p(F-C) < F$ whenever $p>0$, hence
\begin{align}
  r_D < y. \label{r2be5-ineq}
\end{align}

\step{3}{Illustration} $F=100$, $P=90$, $p=5\%$, $C=40$:
\begin{align}
  y &= \frac{100}{90}-1 = 11.11\%, \nonumber\\
  \E[\text{payoff}] &= (0.95)(100)+(0.05)(40) = 95+2 = 97, \qquad r_D = \frac{97}{90}-1 = 7.78\%. \nonumber
\end{align}
\answer{\unskip} The yield is computed from the promised payments, which are the most the holder can receive; with a positive
probability of default and partial recovery the expected payments are lower, so the expected return is below the yield.
The yield spread over Treasuries pays for expected default losses \emph{and} a risk premium; only the latter adds to
the expected return.

\intuition Using the yield as $r_D$ in a WACC overstates the cost of debt, and the error is largest for junk debt,
exactly when default is most likely (Berk Ch.~13 cites KB Homes: 8.6\% promised, about 7.3\% expected).
\pitfalls
\begin{itemize}
  \item Treating a high yield as a high expected return: a bond close to default can promise 20\% and be expected to earn far less.
  \item Using the coupon rate as the cost of debt; it is historical and is neither the yield nor the expected return.
\end{itemize}
\refbib{Berk \S6.5 ``Corporate Bond Yields'', p.~166; Ch.~13 box ``Using the Coupon Rate as the Cost of Debt'', p.~387; Ricca 2A.10 (slides 103--105).}

% =============================================================================
\exsec{6. Credit spreads on one-year zeros --- Berk 6.24, p.~174}
%% BEGIN-STATEMENT 6
\stmtbegin
\textbf{Source:} Berk, Ch. 6, Exercise 6.24, p.~174 (PDF 205) · Answer in book: no

The following table summarizes the yields to maturity on several one-year, zero-coupon securities:

{\def\LTcaptype{none} % do not increment counter
\begin{longtable}[]{@{}ll@{}}
\toprule\noalign{}
Security & Yield (\%) \\
\midrule\noalign{}
\endhead
\bottomrule\noalign{}
\endlastfoot
Treasury & 3.1 \\
AAA corporate & 3.2 \\
BBB corporate & 4.2 \\
B corporate & 4.9 \\
\end{longtable}
}

\begin{enumerate}
\def\labelenumi{\alph{enumi}.}
\item
  What is the price (expressed as a percentage of the face value) of a one-year, zero-coupon
  corporate bond with an AAA rating?
\item
  What is the credit spread on AAA-rated corporate bonds?
\item
  What is the credit spread on B-rated corporate bonds?
\item
  How does the credit spread change with the bond rating? Why?
\end{enumerate}
\stmtend
%% END-STATEMENT

\asked (a) The price, in \% of face, of a one-year AAA zero yielding 3.2\%. (b), (c) The credit spreads of AAA and B
bonds over the Treasury. (d) How spreads move with the rating, and why.

\approach One-year zeros: price $=100/(1+y)$. The credit (default) spread is the difference between a corporate yield and
the Treasury yield of the same maturity.

\notation
\begin{itemize}
  \item $y^{T}=3.1\%$, $y^{AAA}=3.2\%$, $y^{BBB}=4.2\%$, $y^{B}=4.9\%$; spread $s^{k}=y^{k}-y^{T}$.
\end{itemize}

\step{1}{Price}
\begin{align}
  P^{AAA} = \frac{100}{1+0.032} = \frac{100}{1.032} = 96.90. \label{r2be6-P}
\end{align}
\answer{(a)} 96.90\% of face value.

\step{2}{Spreads}
\begin{align}
  s^{AAA} = 3.2-3.1 = 0.1\%,\qquad s^{BBB} = 4.2-3.1 = 1.1\%,\qquad s^{B} = 4.9-3.1 = 1.8\%. \label{r2be6-s}
\end{align}
\answer{(b)} 0.1\% (10 basis points). \answer{(c)} 1.8\% (180 basis points).

\step{3}{Pattern} $0.1<1.1<1.8$: spreads rise as the rating falls.
\answer{(d)} The spread widens as the rating worsens, because lower-rated issuers have higher default probabilities
(and lower expected recoveries): their promised yields must cover larger expected losses and a larger risk premium.

\intuition The spread is the market price of credit risk. By Exercise~5, a B bond's 4.9\% is a \emph{promise}; its
expected return is lower.
\pitfalls
\begin{itemize}
  \item Quoting spreads relative to AAA rather than to the Treasury.
  \item Reading the B spread as the expected extra return; part of it only compensates expected default losses.
\end{itemize}
\refbib{Berk \S6.5 ``Corporate Bonds'', ``Corporate Bond Yields'' and ``Bond Ratings'', pp.~164--167.}

% =============================================================================
\exsec{7. A bond's price over time at a constant yield --- Berk 6.16--6.18, p.~173}
%% BEGIN-STATEMENT 7
\stmtbegin
\textbf{Source:} Berk, Ch. 6, Exercises 6.16--6.18 (shared setup), p.~173 (PDF 204) · Answer in
book: no

Suppose that General Motors Acceptance Corporation issued a bond with ten years until maturity, a
face value of \$1000, and a coupon rate of 7\% (annual payments). The yield to maturity on this bond
when it was issued was 6\%. Use this information for Problems 16--18.

\textbf{16.} What was the price of this bond when it was issued?

\textbf{17.} Assuming the yield to maturity remains constant, what is the price of the bond
immediately before it makes its first coupon payment?

\textbf{18.} Assuming the yield to maturity remains constant, what is the price of the bond
immediately after it makes its first coupon payment?
\stmtend
%% END-STATEMENT

\asked (16) The issue price of a 10-year, 7\% annual-coupon \$1,000 bond at a 6\% yield. (17) Its price one year later just
before the first coupon, and (18) just after it, if the yield stays at 6\%.

\approach Formula \eqref{r2be0-bond}. At a constant yield, the value of any bond grows at the yield between coupon dates
(it is a present value rolling forward), and drops by the coupon when it is paid. The price ``after'' is just the price of
a 9-year bond.

\notation
\begin{itemize}
  \item $C=70$, $F=1{,}000$, $y=6\%$; $P_0$: issue price; $P_1^-$, $P_1^+$: prices just before and just after the first coupon.
\end{itemize}

\step{1}{(16) Issue price} The factors at 6\% for 10 years are
\begin{align}
  A(0.06,10) = \frac{1-(1.06)^{-10}}{0.06} = \frac{1-0.558395}{0.06} = 7.360087,\qquad \frac{1}{(1.06)^{10}} = 0.558395, \label{r2be7-f10}
\end{align}
so
\begin{align}
  P_0 = (70)(7.360087) + (1{,}000)(0.558395) = 515.21 + 558.39 = 1{,}073.60. \label{r2be7-P0}
\end{align}
\answer{(16)} \$1,073.60 (a premium bond: coupon 7\% $>$ yield 6\%).

\step{2}{(17) Just before the coupon} One year on, the same cash flows are one year closer; their present value is
$P_0$ carried forward at 6\%:
\begin{align}
  P_1^- = P_0(1+y) = (1{,}073.60)(1.06) = 1{,}138.02. \label{r2be7-Pm}
\end{align}
\answer{(17)} \$1,138.02.

\step{3}{(18) Just after the coupon} The holder of the bond now owns a 9-year bond. At 6\% for 9 years:
\begin{align}
  A(0.06,9) &= \frac{1-0.591898}{0.06} = 6.801692,\qquad \frac{1}{(1.06)^9} = 0.591898, \nonumber\\
  P_1^+ &= (70)(6.801692) + (1{,}000)(0.591898) = 476.12 + 591.90 = 1{,}068.02. \label{r2be7-Pp}
\end{align}
Check: $P_1^- = 70 + P_1^+ = 70 + 1{,}068.02 = 1{,}138.02$, as in \eqref{r2be7-Pm}.
\answer{(18)} \$1,068.02.

\intuition Over the year the holder earns $0.06\times1{,}073.60=64.42$: a coupon of 70 less a capital loss of
$1{,}073.60-1{,}068.02=5.58$. A premium bond loses value as it ages (it is ``pulled to par''), but its return stays at the
yield. The saw-tooth path (rise, drop by the coupon) is how prices move with time alone.
\pitfalls
\begin{itemize}
  \item Answering (17) and (18) with the same number: the coupon is paid between the two instants.
  \item Thinking the price stays at \$1,073.60 because the yield did not change; with constant yield the \emph{return}, not the price, is constant.
\end{itemize}
\refbib{Berk \S6.4 ``Time and Bond Prices'', pp.~158--160.}

% =============================================================================
\exsec{8. Pricing a coupon bond with spot rates --- Berk 6.10, p.~173}
%% BEGIN-STATEMENT 8
\stmtbegin
\textbf{Source:} Berk, Ch. 6, Exercise 6.10, p.~173 (PDF 204) · Answer in book: no

Assume the current Treasury yield curve shows that the spot rates for 6 months, 1 year, and
\(1\tfrac{1}{2}\) years are 1\%, 1.1\%, and 1.3\%, all quoted as semiannually compounded APRs. What
is the price of a \$1000 par, 4\% coupon bond maturing in \(1\tfrac{1}{2}\) years (the next coupon
is exactly 6 months from now)?

\paragraph{Context needed to solve it}

\emph{From Berk Ch. 6, Appendix B ``The Yield Curve and the Law of One Price'', p.~179 (PDF 210) ---
Eq. 6.4:}

{[}\ldots{]} the price of a coupon bond must equal the present value of its coupon payments and face
value discounted at the competitive market interest rates (see Eq. 5.7 in Chapter 5):

\[P = PV(\text{Bond Cash Flows}) = \frac{CPN}{1+YTM_1} + \frac{CPN}{(1+YTM_2)^2} + \cdots + \frac{CPN + FV}{(1+YTM_n)^n} \qquad (6.4)\]

where \(CPN\) is the bond coupon payment, \(YTM_n\) is the yield to maturity of a zero-coupon bond
that matures at the same time as the \(n\)th coupon payment, and \(FV\) is the face value of the
bond.
\stmtend
%% END-STATEMENT

\asked The price of a \$1,000, 4\% semiannual-coupon Treasury maturing in 1.5 years, given spot rates (semiannual APRs)
of 1\% (0.5 years), 1.1\% (1 year) and 1.3\% (1.5 years).

\approach Block 2A.1: by the law of one price, a coupon bond is a portfolio of zeros, one per cash flow, so each cash flow
is discounted at the spot rate of its own date (Berk Eq.~6.4, equation \eqref{r2be0-spot}). Using a single rate would mix
three different prices of time.

\notation
\begin{itemize}
  \item Cash flows: $CF_1=CF_2=0.04\times1{,}000/2=20$, $CF_3=20+1{,}000=1{,}020$ at half-years 1, 2, 3.
  \item Spot rates per half-year: $r_1=1\%/2=0.5\%$, $r_2=1.1\%/2=0.55\%$, $r_3=1.3\%/2=0.65\%$.
\end{itemize}

\step{1}{Compounding factors}
\begin{align}
  (1.0055)^2 = 1.011030,\qquad (1.0065)^3 = 1.019627. \label{r2be8-f}
\end{align}

\step{2}{Discount each cash flow at its own spot rate}
\begin{align}
  P &= \frac{20}{1.005} + \frac{20}{(1.0055)^2} + \frac{1{,}020}{(1.0065)^3}
     = \frac{20}{1.005} + \frac{20}{1.011030} + \frac{1{,}020}{1.019627} \nonumber\\
    &= 19.90 + 19.78 + 1{,}000.37 = 1{,}040.05. \label{r2be8-P}
\end{align}
\answer{\unskip} $P=\$1{,}040.05$.

\step{3}{The bond's own yield (for intuition)} Solving $1{,}040.05 = 20/(1+y)+20/(1+y)^2+1{,}020/(1+y)^3$ gives $2y=1.295\%$
APR: between the 1-year and 1.5-year spot rates and close to the latter, because most of the value comes from the last cash flow.

\intuition The yield to maturity of a coupon bond is a (value-weighted) average of the spot rates of its dates. With an
upward-sloping curve it lies just below the spot rate of its maturity.
\pitfalls
\begin{itemize}
  \item Discounting everything at 1.3\%: gives a slightly lower price because the early coupons are over-discounted.
  \item Using the annual APR as a per-period rate, or compounding the 1.5-year rate for 1.5 rather than 3 half-years.
\end{itemize}
\refbib{Berk Ch.~6 Appendix B ``Valuing a Coupon Bond Using Zero-Coupon Yields'', pp.~178--181; Ricca 2A.1 (slides 2--7).}

% =============================================================================
\exsec{9. Coupon rate for a new issue sold at par --- Ross 7.19, p.~233}
%% BEGIN-STATEMENT 9
\stmtbegin
\textbf{Source:} Ross, Cap. 7, Exercise 7.19, p.~233 (PDF 257) · Answer in book: no

19. Retornos de títulos de dívida {[}OA2{]} A Gomes \& Cia. quer emitir novas obrigações de 20 anos
para alguns projetos de expansão extremamente necessários. A empresa, no momento, tem obrigações
vendidas por \$930, com cupom de 8\%, que fazem pagamentos semestrais e têm vencimento em 20 anos.
Qual taxa de cupom deve ser definida pela empresa para seus novos títulos de dívida se ela quiser
que sejam vendidos ao par?

\paragraph{Context needed to solve it}

\emph{From Ross §7.1, p.~195 (PDF 219) --- the bond value formula, with face value \(F\), coupon
\(C\) per period, \(t\) periods to maturity and return \(r\) per period:}

\[\text{Valor do título de dívida} = C \times [1 - 1/(1+r)^t]/r + F/(1+r)^t\]

\emph{From Ross Example 7.1, p.~195 (PDF 219):} Os retornos desses títulos são cotados em taxa
nominal anual; a taxa cotada é igual à taxa efetiva do período multiplicada pelo número de períodos.
\stmtend
%% END-STATEMENT

\asked The coupon rate that makes new 20-year semiannual bonds sell at par, when the firm's existing 20-year, 8\%
semiannual bonds sell for \$930.

\approach A bond sells at par exactly when its coupon rate equals its yield to maturity. The market's required return on
20-year debt of this issuer is the yield of the existing bonds, so: solve for that yield, then set the new coupon equal to it.
This is the cost of new debt; the old 8\% coupon is history.

\notation
\begin{itemize}
  \item Existing bond: $C=40$ per half-year, $F=1{,}000$, $n=40$, $P=930$; $y$: yield per half-year.
\end{itemize}

\step{1}{Equation for the yield} From \eqref{r2be0-bond},
\begin{align}
  930 = 40\cdot\frac{1-(1+y)^{-40}}{y} + \frac{1{,}000}{(1+y)^{40}}. \label{r2be9-eq}
\end{align}
The bond is at a discount, so $y>4\%$ per half-year.

\step{2}{Bracket and solve} Trial values: at $y=4.5\%$ (9\% APR) the right side is 907.99; at $y=4.25\%$ (8.5\% APR) it is
952.31. The root lies between and, by iteration (or a financial calculator), is $y=4.3736\%$. Check:
\begin{align}
  A(0.043736,40) = 18.738512,\qquad \frac{1}{(1.043736)^{40}} = 0.180460, \nonumber\\
  (40)(18.738512) + (1{,}000)(0.180460) = 749.54 + 180.46 = 930.00. \label{r2be9-chk}
\end{align}

\step{3}{Annual quote} The yield is quoted as an APR:
\begin{align}
  y^{\text{APR}} = 2\times4.3736\% = 8.75\%,\qquad \text{EAR} = (1.043736)^2-1 = 8.94\%. \label{r2be9-apr}
\end{align}
\answer{\unskip} Set the coupon rate at 8.75\% per year (4.375\% per half-year); then the new bonds sell at par.

\intuition The cost of new debt is the yield the market currently demands, not the coupon promised in the past (block 2A.10:
$r_D$ is not the coupon). The old bonds trade at a discount precisely because 8\% is below what the market now asks.
\pitfalls
\begin{itemize}
  \item Answering 8\% (the old coupon).
  \item Reporting the EAR 8.94\% as the coupon rate; coupon rates, like yields, are quoted as APRs.
\end{itemize}
\refbib{Ross \S7.1 ``Valores e retornos dos t\'itulos de d\'ivida'', pp.~192--199.}

% =============================================================================
\exsec{10. Ratings, prices and issue size --- Berk 6.26, p.~175}
%% BEGIN-STATEMENT 10
\stmtbegin
\textbf{Source:} Berk, Ch. 6, Exercise 6.26, p.~175 (PDF 206) · Answer in book: no

HMK Enterprises would like to raise \$10 million to invest in capital expenditures. The company
plans to issue five-year bonds with a face value of \$1000 and a coupon rate of 6.5\% (annual
payments). The following table summarizes the yield to maturity for five-year (annual-pay) coupon
corporate bonds of various ratings:

{\def\LTcaptype{none} % do not increment counter
\begin{longtable}[]{@{}llllll@{}}
\toprule\noalign{}
Rating & AAA & AA & A & BBB & BB \\
\midrule\noalign{}
\endhead
\bottomrule\noalign{}
\endlastfoot
YTM & 6.20\% & 6.30\% & 6.50\% & 6.90\% & 7.50\% \\
\end{longtable}
}

\begin{enumerate}
\def\labelenumi{\alph{enumi}.}
\item
  Assuming the bonds will be rated AA, what will the price of the bonds be?
\item
  How much of the total principal amount of these bonds must HMK issue to raise \$10 million today,
  assuming the bonds are AA rated? (Because HMK cannot issue a fraction of a bond, assume that all
  fractions are rounded to the nearest whole number.)
\item
  What must the rating of the bonds be for them to sell at par?
\item
  Suppose that when the bonds are issued, the price of each bond is \$959.54. What is the likely
  rating of the bonds? Are they junk bonds?
\end{enumerate}

\paragraph{Context needed to solve it}

\emph{From Berk §6.5 ``Corporate Bonds'', ``Bond Ratings'', p.~166 (PDF 197):}

Bonds in the top four categories are often referred to as investment-grade bonds because of their
low default risk. Bonds in the bottom five categories are often called speculative bonds, junk
bonds, or high-yield bonds because their likelihood of default is high and so they promise higher
yields.

\emph{Table 6.6, p.~167 (PDF 198), rating classes (Moody's / Standard \& Poor's):} Investment Grade
Debt --- Aaa/AAA, Aa/AA, A/A, Baa/BBB; Speculative Bonds (``Junk Bonds'') --- Ba/BB, B/B, Caa/CCC,
Ca/CC, C/C, D.
\stmtend
%% END-STATEMENT

\asked (a) The price of a 5-year, 6.5\% annual-coupon \$1,000 bond if rated AA. (b) The total face value HMK must issue to
raise \$10 million at that price. (c) The rating at which the bond sells at par. (d) The rating implied by a price of \$959.54,
and whether that is junk.

\approach The rating maps to a yield (the table), the yield to a price via \eqref{r2be0-bond}. A bond sells at par when its
yield equals its coupon; and a price can be inverted to a yield and hence to a rating.

\notation
\begin{itemize}
  \item $C=65$, $F=1{,}000$, $n=5$; $y^{k}$: yield for rating $k$ (6.20, 6.30, 6.50, 6.90, 7.50\% for AAA to BB).
\end{itemize}

\step{1}{(a) Price at AA ($y=6.3\%$)}
\begin{align}
  A(0.063,5) &= \frac{1-(1.063)^{-5}}{0.063} = \frac{1-0.736773}{0.063} = 4.178207, \nonumber\\
  P^{AA} &= (65)(4.178207) + (1{,}000)(0.736773) = 271.58 + 736.77 = 1{,}008.36. \label{r2be10-P}
\end{align}
\answer{(a)} \$1,008.36.

\step{2}{(b) Number of bonds} Proceeds per bond are \$1,008.36, so
\begin{align}
  \frac{10{,}000{,}000}{1{,}008.36} = 9{,}917.13 \;\to\; 9{,}917 \text{ bonds},\qquad 9{,}917\times1{,}000 = \$9{,}917{,}000 \text{ of face value}. \label{r2be10-n}
\end{align}
(They raise $9{,}917\times1{,}008.36=\$9{,}999{,}870.56$, within rounding of \$10 million.)
\answer{(b)} \$9.917 million of principal (9,917 bonds).

\step{3}{(c) Par} At $y=6.5\%$ the coupon equals the yield, so each coupon exactly pays the required return and $P=1{,}000$.
\answer{(c)} A rating (yield 6.50\%).

\step{4}{(d) Implied rating} The price is below par, so the yield exceeds 6.5\%. Try BB, $y=7.5\%$:
\begin{align}
  A(0.075,5) &= \frac{1-0.696559}{0.075} = 4.045885,\qquad \frac{1}{(1.075)^5}=0.696559, \nonumber\\
  P &= (65)(4.045885) + (1{,}000)(0.696559) = 262.98 + 696.56 = 959.54. \label{r2be10-d}
\end{align}
The price matches, so the yield is 7.50\%.
\answer{(d)} BB. Yes, they are junk bonds: BB is below the four investment-grade classes (AAA, AA, A, BBB).

\intuition Each notch down the rating scale raises the yield and lowers what the firm raises per bond. For HMK, the difference
between AA and BB is $1{,}008.36-959.54=48.82$ per bond, almost 5\% of the issue.
\pitfalls
\begin{itemize}
  \item In (b), dividing \$10 million by \$1,000 face (10,000 bonds) instead of by the price.
  \item In (d), calling BBB junk: BBB is the lowest investment grade.
\end{itemize}
\refbib{Berk \S6.5 ``Bond Ratings'', Table 6.6 and ``Corporate Yield Curves'', pp.~166--167.}

% =============================================================================
\exsec{11. Interest-rate sensitivity: maturity and coupon --- Berk 6.21--6.22, p.~174}
%% BEGIN-STATEMENT 11
\stmtbegin
\textbf{Source:} Berk, Ch. 6, Exercises 6.21--6.22 (shared table), p.~174 (PDF 205) · Answer in
book: no

Consider the following bonds for Questions 21 and 22:

{\def\LTcaptype{none} % do not increment counter
\begin{longtable}[]{@{}lll@{}}
\toprule\noalign{}
Bond & Coupon Rate (annual payments) & Maturity (years) \\
\midrule\noalign{}
\endhead
\bottomrule\noalign{}
\endlastfoot
A & 0\% & 15 \\
B & 0\% & 10 \\
C & 4\% & 15 \\
D & 8\% & 10 \\
\end{longtable}
}

\textbf{21.} What is the percentage change in the price of each bond if its yield to maturity falls
from 6\% to 5\%?

\textbf{22.} Which of the bonds A--D is most sensitive to a 1\% drop in interest rates from 6\% to
5\% and why? Which bond is least sensitive? Provide an intuition explanation for your answer.

\paragraph{Context needed to solve it}

\emph{Face value:} the statement gives none; Berk's problems in this section use \$1000 bonds
(Problems 13--19), and percentage changes do not depend on the face value.
\stmtend
%% END-STATEMENT

\asked (21) The percentage price change of each of four bonds (two zeros, two coupon bonds; 10 and 15 years) when the yield
falls from 6\% to 5\%. (22) Which is most and which least sensitive, and why.

\approach Price each bond at both yields with \eqref{r2be0-bond} ($F=1{,}000$; percentage changes do not depend on $F$).
Sensitivity grows with the time at which value is received: longer maturity and lower coupon both push value later. The
summary measure is Macaulay duration, the value-weighted average time of the cash flows.

\notation
\begin{itemize}
  \item $P(y)$: price at yield $y$; $\Delta\% = P(5\%)/P(6\%)-1$; $D=\sum_t t\,PV(CF_t)/P$: Macaulay duration at 6\%.
\end{itemize}

\step{1}{Factors} $A(0.06,15)=9.712249$, $A(0.05,15)=10.379658$, $A(0.06,10)=7.360087$, $A(0.05,10)=7.721735$;
$(1.06)^{-15}=0.417265$, $(1.05)^{-15}=0.481017$, $(1.06)^{-10}=0.558395$, $(1.05)^{-10}=0.613913$.

\step{2}{Prices and changes} For example, bond C:
\begin{align}
  P(6\%) &= (40)(9.712249)+(1{,}000)(0.417265) = 388.49+417.27 = 805.76, \nonumber\\
  P(5\%) &= (40)(10.379658)+(1{,}000)(0.481017) = 415.19+481.02 = 896.20, \nonumber\\
  \Delta\% &= 896.20/805.76-1 = 11.23\%. \label{r2be11-C}
\end{align}
In full:
\begin{center}\small\setlength{\tabcolsep}{4pt}
\begin{tabular}{@{}lccccc@{}}\toprule
Bond & Coupon, mat. & $P(6\%)$ & $P(5\%)$ & $\Delta\%$ & $D$ (years)\\\midrule
A & 0\%, 15 & $1{,}000\times0.417265=417.27$ & $1{,}000\times0.481017=481.02$ & $+15.28\%$ & 15.00\\
B & 0\%, 10 & $1{,}000\times0.558395=558.39$ & $1{,}000\times0.613913=613.91$ & $+9.94\%$ & 10.00\\
C & 4\%, 15 & $388.49+417.27=805.76$ & $415.19+481.02=896.20$ & $+11.23\%$ & 11.11\\
D & 8\%, 10 & $588.81+558.39=1{,}147.20$ & $617.74+613.91=1{,}231.65$ & $+7.36\%$ & 7.45\\\bottomrule
\end{tabular}
\end{center}
\answer{(21)} A $+15.28\%$, B $+9.94\%$, C $+11.23\%$, D $+7.36\%$.

\step{3}{Ranking} $15.28>11.23>9.94>7.36$: A $>$ C $>$ B $>$ D, which is the order of the durations $15>11.11>10>7.45$.
\answer{(22)} Most sensitive: A (15-year zero); least sensitive: D (10-year, 8\% coupon). A receives all its value at year 15,
where a change in the discount rate is compounded 15 times; D receives large coupons early, which are barely affected by the rate.

\intuition The relevant ``maturity'' for interest-rate risk is duration, not the final date: C matures in 15 years like A but its
coupons pull its duration to 11.1 years, so it is less sensitive than A and (because 11.1 $>$ 10) more sensitive than the 10-year zero B.
\pitfalls
\begin{itemize}
  \item Ranking by maturity alone (would put C level with A).
  \item Computing the change relative to the new price instead of the old one.
\end{itemize}
\refbib{Berk \S6.4 ``Interest Rate Changes and Bond Prices'', pp.~156--158 (duration, footnote 4, p.~161).}

% =============================================================================
\exsec{12. A bond with deferred, stepped payments --- Ross 7.31 (DESAFIO), p.~234}
%% BEGIN-STATEMENT 12
\stmtbegin
\textbf{Source:} Ross, Cap. 7, Exercise 7.31 (DESAFIO), p.~234 (PDF 258) · Answer in book: no

31. Avaliação de obrigações {[}OA2{]} A Corporação Pereira tem duas obrigações diferentes em
circulação. O título M tem valor de face de \$20.000 e vencimento em 20 anos. Ele não faz pagamentos
nos primeiros seis anos; em seguida, paga \$1.100 a cada seis meses durante os oito anos
subsequentes; e, por último, paga \$1.400 a cada seis meses nos seis anos restantes. O título N
também tem valor de face de \$20.000 e vencimento em 20 anos; ele não tem cupom. Se o retorno
exigido de ambas as obrigações for de 7\% capitalizados semestralmente, qual é o preço atual do
título M? E do título N?
\stmtend
%% END-STATEMENT

\asked The price of bond M (\$20,000 face, 20 years; nothing for 6 years, then \$1,100 every six months for 8 years, then
\$1,400 every six months for 6 years, plus the face value at maturity) and of zero-coupon bond N (\$20,000 face, 20 years),
both at 7\% compounded semiannually.

\approach Value each block of level payments as an annuity at the date just before its first payment, then discount that
value back to today (a deferred annuity); add the face value discounted 40 half-years. Bond N is that last term alone.

\notation
\begin{itemize}
  \item $r=7\%/2=3.5\%$ per half-year; periods $1$--$12$: no payment; $13$--$28$: \$1,100 (16 payments); $29$--$40$: \$1,400 (12 payments); $F=20{,}000$ at $40$.
\end{itemize}

\step{1}{Factors}
\begin{align}
  A(0.035,16) &= 12.094117,\quad A(0.035,12) = 9.663334, \nonumber\\
  (1.035)^{12} &= 1.511069,\quad (1.035)^{28}=2.620172,\quad (1.035)^{40}=3.959260. \label{r2be12-f}
\end{align}

\step{2}{First block (periods 13--28)} Its value at period 12 is an ordinary 16-period annuity; we discount it 12 periods:
\begin{align}
  PV_1 = \frac{1{,}100\times A(0.035,16)}{(1.035)^{12}} = \frac{(1{,}100)(12.094117)}{1.511069} = \frac{13{,}303.53}{1.511069} = 8{,}804.05. \label{r2be12-b1}
\end{align}

\step{3}{Second block (periods 29--40)} Its value at period 28, discounted 28 periods:
\begin{align}
  PV_2 = \frac{1{,}400\times A(0.035,12)}{(1.035)^{28}} = \frac{(1{,}400)(9.663334)}{2.620172} = \frac{13{,}528.67}{2.620172} = 5{,}163.27. \label{r2be12-b2}
\end{align}

\step{4}{Face value}
\begin{align}
  PV_3 = \frac{20{,}000}{(1.035)^{40}} = \frac{20{,}000}{3.959260} = 5{,}051.45. \label{r2be12-F}
\end{align}

\step{5}{Prices}
\begin{align}
  P_M = 8{,}804.05 + 5{,}163.27 + 5{,}051.45 = 19{,}018.78,\qquad P_N = 5{,}051.45. \label{r2be12-P}
\end{align}
\answer{\unskip} $P_M=\$19{,}018.78$; $P_N=\$5{,}051.45$.

\intuition Both bonds promise the same \$20,000 at year 20; M is worth \$13,967 more because of the coupons, even though none is paid for six
years. Deferral is handled by a two-stage discount: annuity value at the start of the block, then a lump sum back to today.
\pitfalls
\begin{itemize}
  \item Discounting the first block's annuity value by 13 periods (its first payment date) instead of 12 (one period before it).
  \item Forgetting the \$20,000 face value of bond M.
  \item Using 7\% per year with annual periods: the payments and the compounding are semiannual.
\end{itemize}
\refbib{Ross \S6.2 (annuities), pp.~155--163, and \S7.1, pp.~192--199.}

% =============================================================================
\exsec{13. Spot and annuity yields from par yields --- Grinblatt \& Titman 10A.3, p.~369}
%% BEGIN-STATEMENT 13
\stmtbegin
\textbf{Source:} Grinblatt \& Titman, Ch. 10, Exercise 10A.3, p.~369 (PDF 382) · Answer in book: no

Compute spot yields and annuity yields for years 1, 2, 3, and 4 if par yields for years 1, 2, 3, and
4 are, respectively, 4.5 percent, 5 percent, 5.25 percent, and 5.25 percent. Assume annual
compounding for all rates and annual payments for all bonds.

\paragraph{Context needed to solve it}

\emph{From Grinblatt \& Titman Appendix 10A, §10A.1 ``Term Structure Varieties'', p.~363 (PDF 376):}

There are a variety of term structures. The par yield curve for the United States represents the
yields and interpolated yields of on-the-run coupon-paying Treasury securities of various maturities
(see Chapter 2). The annuity term structure represents the yields to maturity of riskless bonds with
level payments. There also is the spot term structure, which represents the yields to maturity of
zero-coupon bonds of various maturities.

\emph{From §10A.2 ``Spot Rates, Annuity Rates, and Par Rates'', ``Using Spot Yields to Compute
Annuity Yields and Par Yields'', p.~364 (PDF 377):}

Computing annuity yields from spot yields is relatively straightforward. The annuity yield for year
\(t\) is found by assuming a cash payment of a fixed amount for years \(1, 2, \ldots, t\)---for
example, \$1 is paid at each of those years. This is like a portfolio of \(t\) zero-coupon bonds,
each with a \$1 face value. Using spot yields, compute the present value of each zero-coupon bond
and sum up the \(t\) present values. If the present value is \(P\), find the internal rate of return
for the cash flows: \(-P, 1, 1, \ldots, 1\). This is the annuity yield-to-maturity for date \(t\).
Repeat this process for \(t = 1\), \(t = 2\), \(t = 3, \ldots\), and \(t = T\).

\emph{From ``Computing Spot Yields from Par Yields and Annuity Yields'', pp.~366--367 (PDF
379--380):}

There are many ways to compute the spot yield curve from the par yield curve. Perhaps the easiest
way to proceed is to move sequentially from the shortest maturity to the longest. With annual
compounding and annual coupons, the par bond with the earliest maturity is a zero-coupon bond. The
spot yield is identical to the par yield in this case. The two-year spot yield is computed from a
long position in the two-year par bond and a short position in the one-year par bond. Because the
one-year par bond is a zero-coupon bond, it is possible to form a portfolio that looks like a
zero-coupon bond with two years' maturity by weighting the two par bonds properly. The yield of this
properly weighted portfolio is the spot yield for two years' maturity. Given the one- and two-year
spot rates, it is possible to form a portfolio consisting of a long position in a three-year par
bond and short positions in one- and two-year zero-coupon bonds that look like a three-year zero
coupon bond. The yield of this bond is the three-year spot rate. Proceed sequentially in this manner
as far out into time as necessary.

\emph{From Example 10A.2, p.~367 (PDF 380), the year-3 step with par yields 9\%, 8\%, 7\% (two-year
spot 7.96\%):} To find the year 3 spot rate, strip the two \$7 coupons at years 1 and 2 from the par
bond. This leaves a bond with a value of

\[\$100 - \frac{\$7}{1.09} - \frac{\$7}{(1.0796)^2} = \$87.5722\]

With a future payoff at year 3 in the amount of \$107, and no future cash flows at other dates, the
stripped bond now has the payoff of a zero-coupon bond with a 3-year spot yield \(r\) that solves

\[\$87.5722 = \frac{\$107}{(1+r)^3}\]

implying \(r = 6.91\) percent.
\stmtend
%% END-STATEMENT

\asked Given par yields of 4.5\%, 5\%, 5.25\%, 5.25\% for years 1--4 (annual coupons and compounding), compute the spot (zero-coupon) yields
and the annuity yields for years 1--4.

\approach Bootstrapping (G\&T \S10A.2): a par bond is worth 100 and pays its par yield as coupon. Starting from year 1 (a par bond
of one year is a zero), strip the earlier coupons using the discount factors already found; what is left is a zero-coupon payment
at the final date whose price gives the next discount factor. Annuity yields then follow from the discount factors: the
annuity's price is the sum of the discount factors and its yield is the IRR.

\notation
\begin{itemize}
  \item $c_t$: par yield (coupon per 100) for maturity $t$; $d_t=1/(1+s_t)^t$: discount factor; $s_t$: spot yield; $A_t=\sum_{k\le t}d_k$: price of a \$1 annuity; $a_t$: annuity yield, solving $A(a_t,t)=A_t$.
\end{itemize}

\step{1}{Bootstrap equation} For the $t$-year par bond, price 100 equals coupons plus final payment discounted with the
spot factors:
\begin{align}
  100 = c_t\sum_{k=1}^{t-1}d_k + (100+c_t)\,d_t \;\Longrightarrow\; d_t = \frac{100-c_t\sum_{k=1}^{t-1}d_k}{100+c_t},\qquad s_t = d_t^{-1/t}-1. \label{r2be13-boot}
\end{align}

\step{2}{Year 1} $d_1 = 100/104.5 = 0.956938$, so $s_1=4.50\%$.

\step{3}{Year 2} Strip the year-1 coupon: $5d_1=5\times0.956938=4.784689$.
\begin{align}
  d_2 = \frac{100-4.784689}{105} = \frac{95.215311}{105} = 0.906812,\qquad s_2 = (0.906812)^{-1/2}-1 = 5.0126\%. \label{r2be13-2}
\end{align}

\step{4}{Year 3} Strip two coupons: $5.25(d_1+d_2)=5.25\times1.863750=9.784689$.
\begin{align}
  d_3 = \frac{100-9.784689}{105.25} = \frac{90.215311}{105.25} = 0.857153,\qquad s_3 = (0.857153)^{-1/3}-1 = 5.2723\%. \label{r2be13-3}
\end{align}

\step{5}{Year 4} Strip three coupons: $5.25(d_1+d_2+d_3)=5.25\times2.720903=14.284740$.
\begin{align}
  d_4 = \frac{100-14.284740}{105.25} = \frac{85.715260}{105.25} = 0.814397,\qquad s_4 = (0.814397)^{-1/4}-1 = 5.2667\%. \label{r2be13-4}
\end{align}
Check: $5.25\times2.720903+105.25\times0.814397 = 14.284740+85.715260=100$.
\answer{(spot)} $s_1=4.50\%$, $s_2=5.01\%$, $s_3=5.27\%$, $s_4=5.27\%$ (5.2723\% and 5.2667\%: the 4-year spot is slightly \emph{below} the 3-year).

\step{6}{Annuity prices} Summing the discount factors:
\begin{align}
  A_1=0.956938,\quad A_2=1.863750,\quad A_3=2.720903,\quad A_4=3.535300. \label{r2be13-A}
\end{align}

\step{7}{Annuity yields} Solve $\sum_{k=1}^{t}(1+a_t)^{-k}=A_t$ for each $t$ (one unknown, monotone, so iteration converges):
\begin{align}
  a_1 = 4.50\%,\quad a_2 = 4.8356\%,\quad a_3 = 5.0460\%,\quad a_4 = 5.1296\%. \label{r2be13-a}
\end{align}
For example $t=2$: $1/1.048356+1/1.048356^2 = 0.953874+0.909876 = 1.863751\approx A_2=1.863750$ (the gap is the rounding of $a_2$).
\answer{(annuity)} $a_1=4.50\%$, $a_2=4.84\%$, $a_3=5.05\%$, $a_4=5.13\%$.

\intuition G\&T's rule: when the curve rises, the spot curve is the steepest, the par curve next and the annuity curve the flattest
(here spots $>$ par $>$ annuity at years 2--4). The annuity puts weight on early, low-rate years, so its yield is lowest; a par bond
puts most weight on the final payment, so its yield is close to, but below, the spot rate. The flat par yield from 3 to 4 years needs a
slightly falling spot rate, because the extra coupon is discounted at the higher 3-year rate.
\pitfalls
\begin{itemize}
  \item Setting spot yields equal to par yields beyond year 1.
  \item Discounting the stripped coupons at the par yield instead of at the spot factors already found.
  \item Averaging spot yields to get the annuity yield; it must be solved as an IRR.
\end{itemize}
\refbib{Grinblatt \& Titman Appendix 10A, \S10A.1--10A.2 and Examples 10A.1--10A.2, pp.~363--367.}

% =============================================================================
\exsec{14. The implied value of a call option --- Ross 7.32 (DESAFIO), p.~234}
%% BEGIN-STATEMENT 14
\stmtbegin
\textbf{Source:} Ross, Cap. 7, Exercise 7.32 (DESAFIO), p.~234 (PDF 258) · Answer in book: no

32. Avaliação da opção de resgate antecipado {[}OA2{]} Considere a seguir os preços das três
emissões do Tesouro em 15 de maio de 2007:

{\def\LTcaptype{none} % do not increment counter
\begin{longtable}[]{@{}llllll@{}}
\toprule\noalign{}
\endhead
\bottomrule\noalign{}
\endlastfoot
6,500 & 13 de maio & 106:10 & 106:12 & −13 & 5,28 \\
8,250 & 13 de maio & 103:14 & 103:16 & − 3 & 5,24 \\
12,000 & 13 de maio & 134:25 & 134:31 & −15 & 5,32 \\
\end{longtable}
}

O título do meio é resgatável em fevereiro de 2008. Qual é o valor implícito da opção de resgate
antecipado? (Dica: existe uma forma de combinar as duas emissões não resgatáveis para criar uma
emissão que tenha o mesmo cupom do título resgatável?)

\paragraph{Context needed to solve it}

\emph{Table transcribed from the PDF image (the MD lost the layout); the book prints it without
column headings. The columns follow the layout of the Treasury quotes in Figure 7.4, explained
below. Editor's note (not in the book): ``13 de maio'' translates ``May 13'' of the English
original, i.e.~maturity in May 2013; what matters for the solution is that the three issues mature
on the same date.}

\emph{From Ross §7.5, ``Relatórios de preços de títulos de dívida'' (Treasury quotes, Figure 7.4),
pp.~218--219 (PDF 242--243):}

A Figura 7.4 mostra uma parte das listagens diárias de notas e obrigações do Tesouro encontradas no
site wsj.com. A entrada que começa com ``2021 Nov 15'' está destacada. Essa informação indica que
aquele título vence em novembro de 2021. A coluna ao lado representa a taxa de cupom, que é de
8,000\%. Todos os títulos do Tesouro norte-americano fazem pagamentos semestrais e têm valor de face
de \$1.000, de modo que esse título pagará \$40 a cada seis meses até o vencimento.

As duas informações seguintes são os preços de oferta de compra e de venda. Em qualquer mercado de
balcão ou em corretoras, o preço de oferta de compra (bid) representa aquilo que um corretor está
disposto a pagar por um título, e o preço de oferta de venda (ask) é aquilo que um corretor está
disposto a aceitar por ele. A diferença entre os dois é chamada de margem (ou spread) entre compra e
venda e representa o lucro do corretor.

Por motivos históricos, os preços dos títulos emitidos pelo Tesouro norte-americano são cotados em
32 avos. Assim, o preço de oferta de venda do título de 8\% de novembro de 2021, 136:29, na verdade,
pode ser traduzido como 136 29/32, ou 136,90625\% do valor de face. Com um valor de face de \$1.000,
isso representa \$1.369,0625. Como os preços são cotados em 32 avos, a menor variação de preço
possível é 1/32. Isso é chamado de tick size (tamanho da variação admitida).

O próximo número cotado é o preço de oferta de venda do dia anterior, medida em ticks (ou seja, em
32 avos), de modo que o preço de oferta de venda da emissão subiu em 5/32 de 1\% ou 0,15625\% do
valor de face em relação ao dia anterior. O último número reportado é o retorno até o vencimento,
com base no preço de oferta de venda.
\stmtend
%% END-STATEMENT

\asked The value of the issuer's call option embedded in the 8.25\% Treasury (callable from February 2008), using two noncallable
Treasuries with the same maturity and coupons of 6.5\% and 12\%.

\approach The law of one price again. A portfolio of the two noncallable bonds with the right weights has an 8.25\% coupon and the
same maturity as the callable bond: it is the callable bond \emph{without} the call. The call's value is the price gap
(Exercise~4: $P^{\text{callable}}=P^{\text{straight}}-V^{\text{call}}$). Prices are in 32nds of 1\% of face (Ross's quote convention).

\notation
\begin{itemize}
  \item $w$: weight on the 6.5\% bond, $1-w$ on the 12\% bond (fractions of \$100 face); ask prices $106{:}12$, $103{:}16$, $134{:}31$; bid prices $106{:}10$, $103{:}14$, $134{:}25$.
\end{itemize}

\step{1}{Match the coupon} Same maturity, so matching coupons matches every cash flow:
\begin{align}
  6.5\,w + 12\,(1-w) = 8.25 \;\Longrightarrow\; 12-5.5\,w = 8.25 \;\Longrightarrow\; w = \frac{3.75}{5.5} = 0.681818,\quad 1-w = 0.318182. \label{r2be14-w}
\end{align}
Check: $0.681818\times6.5+0.318182\times12 = 4.431818+3.818182 = 8.25$. Face values match too: $w+(1-w)=1$.

\step{2}{Convert the quotes} $106{:}12 = 106+12/32 = 106.375$; $103{:}16 = 103.5$; $134{:}31 = 134+31/32 = 134.96875$ (\% of face).

\step{3}{Price of the synthetic noncallable 8.25\% bond (ask)}
\begin{align}
  P^{\text{straight}} = (0.681818)(106.375) + (0.318182)(134.96875) = 72.5284 + 42.9446 = 115.4730. \label{r2be14-syn}
\end{align}

\step{4}{Call value}
\begin{align}
  V^{\text{call}} = P^{\text{straight}} - P^{\text{callable}} = 115.4730 - 103.5 = 11.9730\% \text{ of face} \;\Rightarrow\; \$119.73 \text{ per } \$1{,}000. \label{r2be14-call}
\end{align}
With bid prices ($106.3125$, $103.4375$, $134.78125$): synthetic $115.3707$, call $= 11.9332\%$, i.e.\ \$119.33.
\answer{\unskip} About \$119.73 per \$1,000 of face value (11.97\% of par) using ask prices; \$119.33 using bid prices.

\intuition The callable bond trades at 103.5 although an otherwise identical noncallable bond is worth 115.5: the Treasury can retire
it at par in less than a year, so holders will not pay much above 100 plus the coupons until then. The option is deep in the money
(market yields of about 5.25\% versus an 8.25\% coupon), which is why it is worth about 12\% of face.
\pitfalls
\begin{itemize}
  \item Reading $106{:}12$ as $106.12$; the digits after the colon are 32nds.
  \item Mixing bid and ask prices across the bonds.
  \item Comparing yields instead of prices: the callable's quoted yield (5.24\%) is not a yield on comparable cash flows.
\end{itemize}
\refbib{Ross \S7.2 ``Op\c{c}\~ao de resgate antecipado'', p.~207; \S7.5 ``Relat\'orios de pre\c{c}os de t\'itulos de d\'ivida'', pp.~218--219.}

% =============================================================================
\exsec{15. A defaultable bond on a binomial tree --- Grinblatt \& Titman 7.3, p.~254}
%% BEGIN-STATEMENT 15
\stmtbegin
\textbf{Source:} Grinblatt \& Titman, Ch. 7, Exercise 7.3, p.~254 (PDF 267) · Answer in book: no

Value a risky corporate bond, assuming that the risk-free interest rate is 4 percent per period,
where a period is defined as six months. The corporate bond has a face value of \$100 payable two
periods from now and pays a 5 percent coupon per period; that is, interest payments of \$5 at the
end of both the first period and the second period.

The corporate bond is a derivative of the assets of the issuing firm. Assume that the assets
generate sufficient cash to pay off the promised coupon one period from now. In particular, the
corporation has set aside a reserve fund of \$5/1.04 per bond to pay off the promised coupon of the
bond one period from now. Two periods from now, there are three possible states. In one of those
states, the assets of the firm are not worth much and the firm defaults, unable to generate a
sufficient amount of cash. Only \$50 of the \$105 promised payment is made on the bond in this
state.

The exhibit below describes the value of the firm's assets per bond (less the amount in the reserve
fund maintained for the intermediate coupon) and the cash payoffs of the bond. The nonreserved
assets of the firm are currently worth \$100 per bond. At the \(U\) and \(D\) nodes, the reserve
fund has been depleted and the remaining assets of the firm per bond are worth \$120 and \$90,
respectively, while they are worth \$300, \$110, and \$50, respectively, in the \(UU\), \(UD\), and
\(DD\) states two periods from now.

\textbf{Paths for the Value of the Firm's Assets Per Bond (Above the Node) Cash Payoffs of a Risky
Bond (Below the Node) in a Two-Period Binomial Tree Diagram} (transcribed from the PDF image; the
tree recombines at \(UD\)):

{\def\LTcaptype{none} % do not increment counter
\begin{longtable}[]{@{}
  >{\raggedright\arraybackslash}p{(\linewidth - 8\tabcolsep) * \real{0.2000}}
  >{\raggedright\arraybackslash}p{(\linewidth - 8\tabcolsep) * \real{0.2000}}
  >{\raggedright\arraybackslash}p{(\linewidth - 8\tabcolsep) * \real{0.2000}}
  >{\raggedright\arraybackslash}p{(\linewidth - 8\tabcolsep) * \real{0.2000}}
  >{\raggedright\arraybackslash}p{(\linewidth - 8\tabcolsep) * \real{0.2000}}@{}}
\toprule\noalign{}
\begin{minipage}[b]{\linewidth}\raggedright
Node
\end{minipage} & \begin{minipage}[b]{\linewidth}\raggedright
Date
\end{minipage} & \begin{minipage}[b]{\linewidth}\raggedright
Reached from
\end{minipage} & \begin{minipage}[b]{\linewidth}\raggedright
Firm's assets per bond (above the node)
\end{minipage} & \begin{minipage}[b]{\linewidth}\raggedright
Bond cash payoff (below the node)
\end{minipage} \\
\midrule\noalign{}
\endhead
\bottomrule\noalign{}
\endlastfoot
0 & 0 & --- & \$100 & --- \\
\(U\) & 1 & 0 (up) & \$120 & \$5 \\
\(D\) & 1 & 0 (down) & \$90 & \$5 \\
\(UU\) & 2 & \(U\) (up) & \$300 & \$105 \\
\(UD\) & 2 & \(U\) (down) or \(D\) (up) & \$110 & \$105 \\
\(DD\) & 2 & \(D\) (down) & \$50 & \$50 \\
\end{longtable}
}

\paragraph{Context needed to solve it}

\emph{From Grinblatt \& Titman §7 ``Risk-Neutral Valuation of Derivatives: The Wall Street
Approach'', pp.~240--241 (PDF 253--254):}

The risk-neutral valuation method is a three-step procedure for valuing derivatives.

\begin{enumerate}
\def\labelenumi{\arabic{enumi}.}
\tightlist
\item
  Identify risk-neutral probabilities that are consistent with investors being risk neutral, given
  the current value of the underlying asset and its possible future values. Risk-neutral
  probabilities are a set of weights applied to the future values of the underlying asset along each
  path. The expected future asset value generated by these probabilities, when discounted at the
  risk-free rate, equals the current value of the underlying asset.
\item
  Multiply each risk-neutral probability by the corresponding future value for the derivative and
  sum the products together.
\item
  Discount the sum of the products in step 2 (the probability-weighted average of the derivative's
  possible future values) at the risk-free rate. This simply means that we divide the result in step
  2 by the sum of 1 plus the risk-free rate.
\end{enumerate}

{[}\ldots{]} Note from Example 7.8 that .87 is the risk-neutral probability, not the actual
probability, of the good state occurring, which remains unspecified.

\emph{From ``A General Formula for Risk-Neutral Probabilities'', p.~242 (PDF 255):}

One determines risk-neutral probabilities from the returns of the underlying asset at each of the
binomial outcomes, and not by the likelihood of each binomial outcome. The risk-neutral
probabilities, \(\pi\), are those probabilities that make the ``expected'' return of an asset equal
the risk-free rate. That is, \(\pi\) must solve

\[\pi u + (1-\pi) d = 1 + r_f\]

where \(r_f\) = risk-free rate, \(u\) = 1 + per period rate of return of the underlying asset at the
up node, \(d\) = 1 + per period rate of return of the underlying asset at the down node. When
rearranged, this says

\[\pi = \frac{1 + r_f - d}{u - d} \qquad (7.4)\]
\stmtend
%% END-STATEMENT

\asked The no-arbitrage value of a two-period (six-month periods) corporate bond paying \$5 at date 1 (secured by a reserve fund) and
promising \$105 at date 2, of which only \$50 is paid in the default state $DD$; risk-free rate 4\% per period; asset values per bond
as in the exhibit. As an extension tied to block 2A.10 we also compute the bond's promised yield and compare it with its expected return.

\approach The bond is a derivative on the firm's assets, so G\&T's risk-neutral valuation applies: find, at each node, the probability
$\pi$ that makes the assets earn the risk-free rate (eq.~7.4), weight the bond's payoffs with these probabilities, and discount at 4\%
backwards through the tree. The date-1 coupon is riskless (reserve fund), so its value is $5/1.04$.

\notation
\begin{itemize}
  \item $r_f=4\%$ per period; asset values $S_0=100$, $S_U=120$, $S_D=90$, $S_{UU}=300$, $S_{UD}=110$, $S_{DD}=50$.
  \item $\pi_0,\pi_U,\pi_D$: risk-neutral up-probabilities at nodes $0$, $U$, $D$; $V_U,V_D$: date-1 values of the date-2 payment; $V_0$: bond value.
\end{itemize}

\step{1}{Risk-neutral probabilities} From $\pi u+(1-\pi)d=1+r_f$, i.e.\ $\pi S_{\text{up}}+(1-\pi)S_{\text{down}}=S(1+r_f)$:
\begin{align}
  \pi_0 &= \frac{100(1.04)-90}{120-90} = \frac{104-90}{30} = \frac{14}{30} = 0.466667, \label{r2be15-p0}\\
  \pi_U &= \frac{120(1.04)-110}{300-110} = \frac{124.8-110}{190} = \frac{14.8}{190} = 0.077895, \label{r2be15-pU}\\
  \pi_D &= \frac{90(1.04)-50}{110-50} = \frac{93.6-50}{60} = \frac{43.6}{60} = 0.726667. \label{r2be15-pD}
\end{align}

\step{2}{Date-1 values of the date-2 payment} At $U$ both successors pay 105, so the payment is riskless there; at $D$ it is
105 (state $UD$) or 50 (state $DD$):
\begin{align}
  V_U &= \frac{\pi_U(105)+(1-\pi_U)(105)}{1.04} = \frac{105}{1.04} = 100.9615, \label{r2be15-VU}\\
  V_D &= \frac{(0.726667)(105)+(0.273333)(50)}{1.04} = \frac{76.30+13.6667}{1.04} = \frac{89.9667}{1.04} = 86.5064. \label{r2be15-VD}
\end{align}

\step{3}{Date-0 value} The date-2 payment is worth, at date 0,
\begin{align}
  \frac{\pi_0V_U+(1-\pi_0)V_D}{1.04} &= \frac{(0.466667)(100.9615)+(0.533333)(86.5064)}{1.04} \nonumber\\
  &= \frac{47.1154+46.1368}{1.04} = \frac{93.2521}{1.04} = 89.6655, \label{r2be15-V2}
\end{align}
and the coupon is worth $5/1.04=4.8077$. Hence
\begin{align}
  V_0 = 4.8077 + 89.6655 = 94.47. \label{r2be15-V0}
\end{align}
(Same result carrying the coupon through the tree: $[\pi_0(V_U+5)+(1-\pi_0)(V_D+5)]/1.04 = 98.2521/1.04 = 94.47$.)
\answer{\unskip} The bond is worth \$94.47 per \$100 of face value.

\step{4}{Cost of default} A default-free twin would be worth $5/1.04+105/1.04^2 = 4.8077+105/1.0816 = 101.89$. The risk-neutral
probability of reaching $DD$ is $(1-\pi_0)(1-\pi_D)=0.533333\times0.273333=0.145778$, and the shortfall there is $105-50=55$:
\begin{align}
  \frac{55\times0.145778}{(1.04)^2} = \frac{8.0178}{1.0816} = 7.4129 = 101.89-94.47. \label{r2be15-loss}
\end{align}

\step{5}{Promised yield vs.\ expected return (block 2A.10)} The promised yield $y$ solves $94.47 = 5/(1+y)+105/(1+y)^2$: $y=8.10\%$
per half-year (16.21\% APR, 16.86\% effective). Under the risk-neutral probabilities, the expected return of the bond is exactly
$r_f=4\%$ per period (e.g.\ at $U$: $105/100.9615-1=4\%$). The promised yield exceeds this expected return by $8.10-4=4.10$
percentage points, all of it compensation for the expected default loss of step 4. Under the true probabilities (not given) the
expected return would be $4\%$ plus a risk premium, still below 8.10\% because the payoffs can only fall short of the promise.

\intuition The bond equals a default-free bond minus the value of the shareholders' option to hand over the assets in $DD$. Its
yield measures the promise; its expected return (the cost of debt $r_D$) is far lower. Using 8.10\% as $r_D$ would double the cost of
this debt.
\pitfalls
\begin{itemize}
  \item Using the same up-probability at every node; the asset returns differ by node ($+20\%/-10\%$ at 0, $+150\%/-8.3\%$ at $U$), so $\pi$ differs.
  \item Discounting the expected payoff under some guessed ``true'' probability at the risk-free rate; only risk-neutral probabilities go with $r_f$.
  \item Taking the promised yield as the bond's expected return.
\end{itemize}
\refbib{Grinblatt \& Titman Ch.~7, ``Risk-Neutral Valuation of Derivatives'', pp.~239--242 (eq.~7.4); Ricca 2A.6--2A.10 (state prices, $r_D$ vs.\ the yield).}

\end{document}
